\section{Conclusion}
We presented \barrel, a Lean~4 backend for Atelier~B that turns industrial B artifacts into ordinary Lean goals over a set-theoretic encoding of the B language.
Our central design choice is to make B's ubiquitous partiality explicit: partial operators are represented as total Lean functions parameterized by proofs of well-definedness which are reified as standard additional goals.
The requirement applies both to expressions translated from the specification and to terms introduced during interactive proof.

The first Atelier~B example closes a contradictory assertion using \(\min\varnothing\); the second does so through a supplied rule whose counterinstance uses only total operators.
In \barrel, the former requires an impossible definedness proof; Lean refutes the latter rule instance using Mathlib.
The leader-election development illustrates how the same proof discipline supports reuse: mathematical lemmas isolate arguments from generated contexts, and inherited WD conditions are justified by explicit applications of earlier components' theorems.
Automation can use these checked facts without becoming an additional trusted inference mechanism.

Lean's kernel checks the resulting arguments, but the external PO generator and the faithfulness of the translation remain trusted.
Future work includes extending the supported B fragment and lemma library, improving the interface to generated hypotheses, and establishing an end-to-end guarantee through a verified translation and proof obligation generator.

\paragraph*{Artifacts.}
A snapshot of \barrel is available at \url{https://anonymous.4open.science/r/BARReL}.
The archive contains the complete Lean~4 implementation of \barrel (B surface embedding, POG reader, encoder, discharger macros, and tactic layer), together with the sample B machines and Lean scripts used in the evaluation.
% ---including the \texttt{MinSearch} refinement chain and the \texttt{DerivFalse} counterexample.
% The artifacts are a standard Lean~4 project and pin their dependencies via \texttt{lean-toolchain} and \texttt{lake-manifest.json}; a typical build is \texttt{lake update} followed by \texttt{lake build}.
% To replay the end-to-end workflow that regenerates proof obligations from B sources, an Atelier~B installation providing \texttt{bxml} and \texttt{pog} is required; which then has to be pointed inside \barrel by setting the directory containing \texttt{bin/} and \texttt{include/} via \texttt{set\_option barrel.atelierb "<path>"}.
